# Workshop: Scoping

Im Scoping-Workshop konkretisieren wir die mit einem Vorhaben verbundenen Ziele und Randbedingungen. Wir definieren, welche (quantitativen) Metriken zur Messung des Erfolgs der bevorstehenden Projektarbeit geeignet sind. Gemeinsam mit dem Projektteam und der Auftraggeberin wollen wir die hinter einem Vorhaben liegenden Annahmen aufdecken und hierbei kritische Annahmen identifizieren.

## Überblick

Der Scoping-Workshop ist der erste in einer Reihe von acht Workshop-Schritten unseres Vorgehensmodells: Er markiert den Beginn der *Verstehen*-Phase des Entwicklungsprozesses. Ziel eines Scoping-Workshops ist es, einen Projektauftrag möglichst umfassend zu verstehen, die Messung des Projekterfolgs zu operationalisieren und herauszufinden, auf welchen Annahmen anvisierte Projektziele basieren.

**Ziel**
- Klärung des Projektauftrags

**Teilnehmende**
- Projektteam, Auftraggeberin

**Hauptergebnisse**
- Problem Statement Map, Proto-Personas

**Ablauf**
- Proto-Problem Statement Map erstellen
- Proto-Personas formulieren
- Proto-Journey formulieren
- Annahmen priorisieren

Ein präzises Verständnis des Projektauftrags ist von entscheidender Bedeutung für dessen zielorientierte Bearbeitung. Allzu oft wird dieser Schritt als trivial angesehen und daher leichtfertig übersprungen. In einer Vielzahl von Projekten konnten wir lernen, wie wichtig die umfassende Klärung des Projektauftrags ist: Gemeinsam mit der Auftraggeberin schärfen wir im Scoping-Workshop den Blick auf die vorliegende Problemstellung.

Im späteren Research-Workshop werden wir Methoden auswählen, mit denen wir die Belastbarkeit der identifizierten Annahmen empirisch prüfen können. Im Synthese-Workshop führen wir anschließend die erhaltenen Ergebnisse zusammen und können auf der so konsolidierten Grundlage Inhalt und Formulierung des Projektauftrags ergänzen und – sofern notwendig – anpassen.

Der Scoping-Workshop verläuft typischerweise in zwei Phasen. In der ersten Phase arbeiten wir den eigentlichen Projektauftrag heraus. In der zweiten Phase identifizieren wir dann die Annahmen, die dem gewonnenen Verständnis des Projektauftrags zugrunde liegen.

Bei kleinen Projekten dauert ein Scoping-Workshop oft lediglich einen halben Tag. Bei sehr ambitionierten Projektvorhaben kann es hingegen lohnenswert sein, den Scoping-Workshop aufzutrennen und entsprechend mehr Zeit zu reservieren: Der erste Teil des Scoping-Workshops kann dann beispielsweise Kick-off-, der zweite Annahmen-Workshop genannt werden.

Begleiten wir nun den Scoping-Workshop in unserem Fallbeispiel zur Leistungserfassung bei 4Service.

## Proto-Problem Statement Map

Im Scoping-Workshop erarbeiten wir erste Ergebnisse und dokumentieren diese in Form modellierter Artefakte, die wir häufig als – physische oder digitale – *Maps* konkretisieren. Maps sind zweidimensionale, aus einzelnen Karten bestehende Strukturen, die sich leicht arrangieren und verändern lassen. Um zu verdeutlichen, dass es sich an dieser frühen Stelle im Projektverlauf noch um vorläufige, nicht validierte Ergebnisse handelt, stellen wir diesen Modellen oder UX-Artefakten den Terminus »Proto« voran, der im Lateinischen »erstmalig« oder »frühest« bedeutet.

Starten wir nun mit dem Scoping-Workshop: Tim hat das Produktmanagement von 4Service eingeladen. Er weiß, wie wichtig ein guter Verlauf des ersten Workshops für die weitere Projektarbeit und damit für den Gesamterfolg eines Projekts ist. Damit sein Team über eine solide Ausgangslage verfügt und alle Mitarbeitenden die mit dem neuen Release von 4Service verbundenen Projektziele umfassend nachvollziehen können, möchte Tim die Erwartungen und Annahmen der Auftraggeberin möglichst detailliert verstehen.

Nachdem alle Beteiligten im Raum Platz genommen haben, bittet Tim die Anwesenden um eine kurze Vorstellungsrunde: Er will die beteiligten Personen und ihre Rollen als *Stakeholder* bekannt machen.

> **Stakeholder:** Stakeholder sind Anspruchsgruppen, die von einer Lösung direkt oder indirekt betroffen sind. Oft haben sie auch ein unmittelbares Interesse an der Lösung.

Die im Workshop zur Verfügung stehende Zeit soll möglichst produktiv genutzt werden: Tim öffnet seinen Moderationskoffer und beginnt damit, Karten an eine vorbereitete Pinnwand zu hängen. Er verwendet zunächst blaue Karten und ordnet sie horizontal in einer Reihe an (Abbildung 4). Mit Überschriften auf den Karten gibt Tim die Struktur eines *Proto-Problem Statement* vor, das er nun erarbeiten möchte. Eine erste Map entsteht.

Die erste Karte beschriftet Tim mit »Nutzende«. Er fragt die Anwesenden: »Zwischen welchen Nutzergruppen sollten wir unterscheiden? Welche Nutzergruppen haben unterschiedliche Anforderungen an die Leistungserfassung?« Tim erhält auf seine Frage eine Reihe von Antworten – nicht alle davon weisen jedoch in die gleiche Richtung: Es finden sich durchaus widersprüchliche Meinungen.

In einem ersten Versuch zur Segmentierung der Nutzenden von 4Service arbeiten die Teilnehmenden schließlich vier unterschiedliche Nutzergruppen heraus. Diese sind:

- Projektleitung
- Teamleitung
- Personalleitung
- Mitarbeitende

> **Problem Statement:** Ein Problem Statement beschreibt die zu bearbeitende Fragestellung. Das Statement definiert, welches Problem konkret zu überwinden ist, wer von dem Problem und dessen Lösung betroffen ist und welche Randbedingungen dabei zu berücksichtigen sind. Im Scoping-Workshop sprechen wir bewusst noch von Proto-Problem Statement, um die Unsicherheit der darin zusammengefassten Annahmen auszudrücken. Im Verlauf der Verstehen-Phase wird das Bild des Problem Statement zunehmend geschärft und gegebenenfalls in Teilen auch revidiert.

Tim hält die Wortmeldungen auf Karten fest und ordnet sie in der Spalte »Nutzende« an: Eine Tabelle entsteht (Abbildung 5).

Als Nächstes fragt Tim: »Welche Probleme möchten wir für die direkten User von 4Service mit Blick auf die Zeiterfassung lösen?« Die Teilnehmenden des Workshops vermuten, dass es viele Reklamationen wegen falscher Rechnungen geben wird. Zum Beispiel dann, wenn eine Anwaltskanzlei, als Kunde von 4Service, für zwei Anwälte, die sich in einer gemeinsamen Sitzung beraten haben, unterschiedliche Zeitaufwände in Rechnung stellt.

Ein weiteres Problem resultiert für Kunden, wenn Arbeiten für einen Klienten zwar ausgeführt, jedoch nicht abgerechnet werden: Der Kanzlei oder einem Unternehmen entstehen Aufwände, denen keine Einnahmen gegenüberstehen. Für dieses Problem lassen sich unterschiedliche Ursachen annehmen: Eine regelmäßige Zeiterfassung könnte für Nutzende zu aufwendig und ineffizient sein und daher vielleicht erst gesammelt am Monatsende durchgeführt werden. Kleinere Arbeiten für einen Klienten »zwischendurch« geraten dann schon mal in Vergessenheit. Da die 4Service AG keine Smartphone-App zur Nutzung von 4Service anbietet, ist eine mobile Erfassung von Arbeitszeiten nicht möglich. Personen, die eine Leistung außerhalb des Büros erbringen, müssen sich die angefallenen Aufwände daher merken oder notieren, bis sie wieder Zugriff auf 4Service haben.

Tim moderiert die Diskussionsrunde zu möglichen Ursachen für ausbleibende Zeiterfassungen, hinterfragt die ausgetauschten Aussagen der Vertreterinnen des Produktmanagements und hält resultierende Erkenntnisse auf Karten fest – sukzessive werden die Spalten des Proto-Problem Statement gefüllt (Abbildung 6).

Die sorgfältige Unterscheidung zwischen Problemen und Ansätzen zu ihrer Lösung ist hilfreich: Wir richten den Fokus zunächst auf das Verständnis der Ausgangslage und die Zielsetzungen eines Projekts.

Tim möchte wissen, ob es unter den Teilnehmenden aus dem Produktmanagement bereits Ideen zu möglichen Lösungsansätzen gibt. Tim weiß, dass es wichtig ist, bestehende Ideen der Auftraggeberin in Erfahrung zu bringen: Zum einen gelingt hierdurch der Einbezug dieser Stakeholder leichter – zum anderen sollen bereits vorliegende Lösungsideen nicht unberücksichtigt gelassen werden. Vorhandene Lösungsideen zu kennen, ist auch deswegen bedeutsam, weil wir dann die dahinter verborgenen Annahmen hinterfragen können.

Mit der Visualisierung der Problembestandteile vor Augen fällt das Nachdenken über mögliche Lösungen im Team leichter. Von Vertreterinnen des Produktmanagements werden auch gleich zwei Ideen für Lösungsansätze genannt: »Einfachere Erfassungsmasken« und ein »mobiler Client« könnten die angesprochenen Probleme lösen. Auch wenn dieser Zeitpunkt – das Team erkundet aktuell noch Annahmen zum Problemraum – für ein informiertes Ausloten des Lösungsraums noch zu früh ist, schafft diese Spalte in der Proto-Problem Statement Map eine Art »Ventil« für jene Stakeholder, die von ihrem Lösungsansatz (oft: allzu) überzeugt sind.

»Wenn wir über Lösungen diskutieren«, sagt Tim, »so brauchen wir Möglichkeiten zur Bewertung ihrer Qualität.« Wie lässt sich der Erfolg vorgeschlagener Lösungen feststellen, wenn wir sie umsetzen? Hierzu bedarf es geeigneter *Metriken* und ihnen zugeordneter Messmethoden.

> **Metrik:** Eine Metrik bildet eine Produkteigenschaft auf ein qualitatives oder quantitatives Gütemaß ab und macht damit das Vorhandensein und/oder die Messung der Ausprägung einer Eigenschaft überprüfbar.

Durch die explizite Definition von Metriken objektiviert das Team die konkrete Zielsetzung und, damit verbunden, welche Ergebnisse als Erfolg bewertet werden. Hierdurch wird ein Vergleich zwischen dem aktuellen Istzustand und dem intendierten Zielzustand nach Projektabschluss möglich. Metriken verlangen nach Messmethoden, ihre Festlegung gibt uns Richtungen vor – oder wie es der Ökonom Peter Drucker formulierte: »Nur was wir messen, können wir auch lenken.«

Die in einem Problem Statement zu bestimmenden Metriken stellen einen Bezug zu den identifizierten Problemen her und beschreiben, wie ein geeigneter Zielzustand aussehen sollte. Zu ihrer Feststellung verlangen Metriken nach einer klaren Operationalisierung und einer hierauf fußenden Zuordnung von Messmethoden. Zur Prüfung der Frage, ob beispielsweise die Leistungserfassung mit 4Service effizient ist, bietet sich die Messung des durchschnittlichen Zeitaufwands zur Erfassung einer Leistung als Operationalisierung der Metrik »Effizienz« an. »Weniger als 10 Minuten Zeit pro Tag für die Leistungserfassung« wird von einem Teilnehmer als quantitatives Ziel für die Metrik »Effizienz« genannt. Als Tim fragt, wie lange die Leistungserfassung denn mit der aktuellen Version von 4Service dauere, wird deutlich, dass hierzu keine Erfahrungswerte vorliegen – lediglich ein unsicheres »vermutlich deutlich länger als 10 Minuten pro Arbeitstag« wird als Annahme geäußert. Die vage Aussage ist ein Hinweis darauf, dass der Istzustand mit Blick auf die Metrik *Effizienz* nicht klar ist. Der Istzustand sollte also überprüft und dann eine Einigung über den Soll- oder Zielzustand getroffen werden. Tim plant dies als Aufgabe für einen nachfolgenden Workshop ein.

Nun fragt Tim nach einem zweiten Problembereich: »Wie steht es um die Verrechenbarkeit der Arbeitszeit von Mitarbeitenden. Dies wurde auch als Problem deklariert.«

Sarah, die Testleiterin im Team, meldet sich: »Diese Größe haben wir als direkten Kennwert in 4Service. Eine Mitarbeiterin kann die Verrechenbarkeit ihrer Arbeitszeit im Leistungserfassungsmodul einsehen und beispielsweise feststellen, dass sie eine durchschnittliche Verrechenbarkeit von 68% ihrer Arbeitszeit hat. Das heißt, 68% der von ihr aufgebrachten Zeit kann Kunden in Rechnung gestellt werden. Sollten tatsächlich Leistungseinträge von Mitarbeitenden vergessen werden, so hat dies negative Auswirkungen auf ihre Verrechenbarkeit. Wenn wir eine genaue Vorstellung über das Ausmaß unprotokollierter Leistungen erfahren und die Ursachen dafür verstehen wollen, müssen wir Nutzerforschung betreiben. In welcher Größenordnung seht ihr die Möglichkeit, durch Verbesserungen in der Leistungserfassung von 4Service eine Erhöhung der tatsächlichen Verrechenbarkeit der Arbeitszeit von Mitarbeitenden zu erreichen?«

Nach einer längeren Diskussion einigt sich das Team darauf, dass durch eine Optimierung des Erfassungsmoduls weniger Zeiten unberücksichtigt blieben und durch eine korrekte und vollständige Protokollierung eine Erhöhung der Verrechenbarkeit in der Größenordnung von etwa 5% realistisch zu erwarten ist. Die Diskussion zeigte jedoch, dass diese Schätzung durch gezielte Forschungstätigkeiten überprüft werden muss.

Ähnliches gilt auch für Kundenbeschwerden wegen falscher Rechnungen, denen Erfassungsfehler zugrunde liegen. Aus einer früheren Analyse liegen Hinweise vor, dass etwa eine von zehn Rechnungen reklamiert wird. Das Team einigt sich darauf, dass diese Anzahl unbedingt halbiert werden sollte. Auch diese Größenordnung wird als Annahme gesehen, die im Verlauf des Projekts zu überprüfen ist.

Tim regt nun die Teilnehmenden dazu an, zu überlegen, ob es neben den im Proto-Problem Statement unter »Nutzende« genannten Rollen weitere Anspruchsgruppen gibt, die im Projektverlauf Berücksichtigung finden sollten. Vor allem die Gruppen »Produktmanagement«, »Supportmitarbeitende« und »Sales/Marketing« werden von den Teilnehmenden des Workshops genannt.

Mit seinen Fragen möchte Tim die *Randbedingungen* ausloten, die von den Beteiligten im Workshop angenommen und deren Einhaltung als erfolgskritisch angesehen werden. Eine präzise Kenntnis der Randbedingungen eines Projekts ist essenziell: Wir müssen die Charakteristika des Problemraums verstehen, damit sie in der Lösungsfindung angemessen Berücksichtigung finden und wir die Vorzüge und Restriktionen etwaiger Lösungsoptionen bewerten können. Als Randbedingung aus technischer Perspektive wird zum einen genannt, dass die Struktur der Datenbank, die dem gegenwärtigen Leistungserfassungsmodul zugrunde liegt, nicht geändert werden soll: Es würden sonst bei Einführung einer neuen Version von 4Service große Aufwände bei der Migration von Daten entstehen. Zum anderen wolle man keine native mobile Lösung entwickeln, sondern an der bisher eingesetzten Webtechnologie festhalten. Tim bespricht mit den Personen im Workshop, offene Fragen zu technischen Implikationen an das Entwicklungsteam zu delegieren, um die Entwicklerinnen auf diese Weise frühzeitig in das Projekt einzubinden.

> **Risiko:** Risiken sind mögliche Probleme, deren Eintreffen aus heutiger Sicht nicht sicher ist. Sie beinhalten eine Eintretenswahrscheinlichkeit und eine Auswirkung. Wir unterscheiden zwischen Projekt- und Produktrisiken. Projektrisiken können sich während der Projektlaufzeit auswirken, wohingegen sich Produktrisiken gegebenenfalls erst nach dem Release eines Produkts manifestieren können.

Im Workshop werden auch *Risiken* des vorliegenden Projekts diskutiert. Dabei identifizieren die Teilnehmenden vor allem zwei Risikofaktoren: ein technisches Risiko und eines, das mit möglicherweise unterschiedlichen Anforderungen der vielfältigen Nutzergruppen verbunden ist.

Das technische Risiko betrifft die Implementierung von 4Service: Die der aktuellen Version zugrunde liegende Webtechnologie könnte aktuellen Anforderungen an ein performantes User Interface nicht mehr gerecht werden. Zur Klärung und Bewertung dieses Risikos sollte ein technischer Prototyp als *Proof of Concept* entwickelt und näher analysiert werden. Als zweiter Risikofaktor könnten sich möglicherweise unterschiedliche Nutzungsanforderungen an ein adäquates Modul zur Leistungserfassung erweisen: Es muss erkundet werden, ob eine gefundene Lösung die verschiedenen relevanten Nutzergruppen angemessen unterstützt.

Aus den Diskussionen im Workshop entsteht schließlich ein umfassendes Proto-Problem Statement, das die erarbeiteten Erkenntnisse zusammenfasst (Abbildung 7).

Natürlich könnten wir ein solches Problem Statement auch verbal formulieren. In der dargestellten Form als Map lässt es sich jedoch im Projektverlauf leicht im Lichte neuer Erkenntnisse anpassen und erweitern. Zudem erlaubt es den Einbezug der unterschiedlichen Perspektiven aller Teilnehmenden und der jeweils relevanten Aspekte. Es ist das dynamisch entwickelte, reichhaltige Ergebnis der bisherigen Teamarbeit.

Tim richtet sich an die Kolleginnen und Kollegen aus dem Produktmanagement: »Herzlichen Dank – mit dem Proto-Problem Statement sind wir bei der Klärung der Ausgangslage ein ganzes Stück vorangekommen. Im nächsten Schritt möchten wir nun herausfinden, wie gesichert unser vermeintliches Wissen ist – wir müssen den Status der Annahmen bestimmen, die sich hinter unserem Proto-Problem Statement verbergen.«

Annahmen sind Vermutungen, die wir anstellen, ohne uns im Hinblick auf ihren Wahrheitsgehalt ganz sicher sein zu können. Oft ist uns der Unterschied zwischen gesichertem Wissen und bloßen, ungeprüften Annahmen gar nicht so klar. Annahmen können implizit sein – wir müssen sie erst als solche »entdecken« und dann die Belastbarkeit ihres Inhalts sorgfältig erschließen. Das Aufspüren impliziter Annahmen kommt dabei nicht selten einer Detektivarbeit gleich. Deswegen sprechen wir von Proto-Problem Statement: Wir möchten die Tatsache herausstellen, dass es sich um ein *angenommenes* Problem Statement handelt, das wir später anhand von Nutzerforschungen in seiner Gültigkeit noch validieren sollten.

Durch das explizite Nachfragen zu Projektzielen und möglichen Lösungsansätzen konnte Tim bereits einige Annahmen offenlegen. Er regt die Leute im Workshop zur kritischen Reflexion der Annahmen im Proto-Problem Statement an: »Stimmt die auf der Karte festgehaltene Aussage tatsächlich? Wieso bin ich mir da eigentlich so sicher?« Unsichere Annahmen können anschließend im Team diskutiert werden: »Könnten wir mit unseren Annahmen zu unterschiedlichen Nutzergruppen falsch liegen?«

Wie die Diskussion in unserem Beispiel zeigt, besteht nicht unter allen Teilnehmenden Einigkeit bezüglich der Korrektheit der festgehaltenen Annahmen. Tim ist davon nicht überrascht, oft werden erst durch ein beharrliches, kritisches Hinterfragen von Annahmen unterschiedliche Perspektiven sichtbar. Annahmen werden erst dann gefährlich, wenn wir ungeprüft unsere Entscheidungen auf sie gründen. Das kollaborative Reflektieren von Annahmen im Team legt offen, welche Forschungsfragen wir beantworten müssen, damit wir ein Problem auf solider Grundlage lösen können. Im vorliegenden Fall sollte beispielsweise untersucht werden, wie viele relevante Nutzergruppen bei der Leistungserfassung tatsächlich zu differenzieren sind oder wie aufwendig die Protokollierung von Arbeitszeiten de facto ist.

## Theorie: Workshop-Grundlagen

### Prinzipien erfolgreicher Workshops

Wir haben Tim gerade bei der Durchführung eines Workshops erlebt. Tim ist ein erfahrener Projektleiter und konnte in vielen – erfolgreichen wie auch weniger erfolgreichen – Workshops Erfahrungen sammeln. Auf dieser Grundlage hat er einige Regeln aufgestellt, die einen positiven Ausgang von Workshops begünstigen: Wir stellen sie in diesem Abschnitt vor.

**Ergebnisorientierte Gestaltung.** Bei effizienten Workshops hat der Leiter oder die Leiterin eines Workshops bereits vor dessen Beginn eine klare Vorstellung von den erreichbaren Ergebnissen – natürlich nicht inhaltlich, wohl aber mit Blick auf deren Struktur. Tim bereitet die Wände des Workshop-Raums entsprechend mit Platzhaltern für die erwarteten Inhalte vor. An der Wand hängen Poster, die beispielsweise mit Karten zu »Aufgaben«, »Annahmen«, »Parkplatz« oder mit »Offene Fragen« überschrieben sind. Die inhaltlichen Arbeitsaufträge für den Workshop sind auf diese Weise für alle Teilnehmenden sichtbar (Abbildung 8). Ein Beispiel für eine solche Vorbereitung haben wir zuvor bereits bei der Formulierung des Proto-Problem Statement gesehen.

**Beiträge sichtbar machen.** Damit ein erreichter Diskussionsstand transparent ist und Beiträge aller Teilnehmenden eines Workshops Berücksichtigung finden, sollten die Beiträge visualisiert werden. Auch wenn Tim als Leiter des Workshops einen Beitrag zunächst für weniger wertvoll hält, so dokumentiert er diesen dennoch. Oft geben auch solche Beiträge im späteren Verlauf Anstöße, die sich durch weiteres Nachfragen zu zielführenden Ergebnissen wandeln. Sind alle Beiträge visualisiert und auf diese Weise dokumentiert, so ist auch das Protokoll bereits erledigt – wir alle kennen die lästige Frage nach der Protokollführung in einem Workshop. Tim hält die erreichten Ergebnisse einfach durch ein Foto der Plakate an den Wänden fest – es bedarf lediglich eines Smartphones und das zentrale Ergebnis ist dokumentiert. Fotos von erreichten Zwischenständen eignen sich zur Protokollierung des groben inhaltlichen Verlaufs eines Workshops. Natürlich sind alle in diesem Buch beschriebenen Methoden auch oder sogar explizit für online durchgeführte Workshops geeignet. Verwendete Onlinetools geben in der Regel keine Platzrestriktionen vor und die Dokumentation ist ohnehin ein inhärenter Vorzug dieser Tools.

**Timeboxed.** Aktivitäten im Workshop sollten zeitlich streng begrenzt werden. Zu schnell verlieren sich auch ansonsten produktive Teams in fruchtlosen Diskussionen. Eng getaktete Aufgaben schaffen ein Klima der Fokussierung und Effizienz. Zugegeben: Zeitbegrenzungen sind anspruchsvoll und anstrengend. Auch ist eine starre zeitliche Taktung meist ungewohnt und erfordert einige Überwindung für die moderierende Person eines Workshops. In der Praxis finden wir auch in Workshops die bekannte 80%-Regel oft bestätigt: In 20% der Zeit erarbeiten wir 80% der Ergebnisse. Oft werden Resultate nicht besser, wenn wir ein Thema unstrukturiert hin und her diskutieren – im Gegenteil: Wir drehen uns im Kreis und stellen bereits Erreichtes ohne neue Erkenntnisse infrage, statt uns konzentriert auf die nächsten Schritte vorzubereiten.

**Logischer Aufbau, klare Struktur.** Ein guter Workshop basiert auf einer geplanten Sequenz logisch aufeinander aufbauender Phasen. Die Abfolge der Phasen bildet einen Problemlösungsprozess ab. Erfolgreiche Problemlösungsprozesse beschäftigen sich zunächst mit dem Verständnis der Ausgangslage und den wesentlichen Eigenschaften eines Problems. Erst dann werden auf der so geschaffenen belastbaren Basis Lösungsideen generiert und bewertet. Wir müssen erst den Problemraum verstehen, bevor wir den Lösungsraum informiert erkunden können. Erfolgversprechende Ideen lassen sich dann leicht von weniger aussichtsreichen Ansätzen trennen.

**Einsam und gemeinsam.** Die vielleicht wichtigste (und gar nicht so geheime) Zutat für gute Workshop-Ergebnisse ist ein abgestimmtes Wechselspiel zwischen gemeinsamen Tätigkeiten und solchen Aktivitäten, die in Einzelarbeit erfolgen. Beispielsweise können Teilnehmende eines *Design Studios* (siehe das Kapitel Ideation-Workshop in diesem Buch) in einem ersten Schritt zuerst eigenständig verschiedene Lösungsvorschläge erarbeiten. In einer zweiten Phase werden die resultierenden individuellen Vorschläge dann einander präsentiert und gemeinsam beurteilt. Danach arbeiten die Teilnehmenden auf der Basis des erhaltenen Feedbacks wiederum alleine an ihren Entwürfen weiter.

Im Verlauf des Fallbeispiels werden wir verschiedene Workshop-Methoden vorstellen, die ein solches Wechselspiel von Einzelarbeit und gemeinsamen Tätigkeiten beinhalten.

**Ausbalanciert.** In guten Workshops kommen alle Personen zu Wort. Auch jene, die sich ansonsten eher zurückhalten: Wir möchten auch ihrer Perspektive Raum geben. Teilnehmende, die üblicherweise bei jeder Gelegenheit viel sprechen, müssen in ihrem Redefluss im Zaum gehalten werden. Der Einsatz geeigneter Methoden hilft, dies sicherzustellen – hier ist die Kompetenz der moderierenden Person eines Workshops gefragt. Statt einfach nach Beiträgen zu einem Thema zu fragen, können wir die Teilnehmenden eines Workshops beispielsweise auffordern, ihre Ideen zunächst für sich auf drei Karten zu schreiben und diese Inhalte anschließend allen vorzustellen. Auf diese Weise involvieren wir alle Anwesenden und machen ihre Beiträge sichtbar: Alle am Workshop Beteiligten tragen gleichermaßen Verantwortung für dessen Ergebnisse.

### Planung von Workshops

Stellen wir uns kurz eine Workshop-Situation vor. Die Moderatorin eines Workshops präsentiert zu Beginn einige gesicherte Fakten zu einer Problemstellung. Danach sammeln die Teilnehmenden verschiedene Lösungsideen, die sie sich wiederum gegenseitig präsentieren. Anschließend erfolgt eine Feedbackrunde: Auf dieser Grundlage wird abgestimmt und die präferierte Lösung ausgewählt.

Wie das Beispiel verdeutlicht, besteht der skizzierte Workshop aus unterschiedlichen, aufeinander folgenden Schritten. Die Planung dieser Schritte und die Auswahl einer zielgerichteten Methode in jeder Phase bildet eine wichtige Grundlage für das Gelingen von Workshops. Nachfolgend stellen wir verschiedene Arten solcher Phasen vor.

**Präsentation.** Oft beginnt ein Workshop mit einer Präsentation. Dabei werden typischerweise Ergebnisse vorgestellt, die zuvor alleine oder durch ein kleines Team erarbeitet wurden und die nun mit der gesamten Gruppe geteilt werden. Präsentationen bauen häufig auf vorbereiteten Folien auf – dies bringt einen großen Nachteil mit sich: Einmal gezeigt, sind die Folien danach nicht mehr zu sehen. Weder die verwendeten Folien noch die Kommentare, die zu diesen geäußert wurden, bleiben für die Teilnehmenden eines Workshops über eine gewisse Zeit präsent. Besser ist es, präsentierte Inhalte im Raum für alle zu visualisieren oder auf einem virtuellen Board zur Verfügung zu stellen. So können alle Personen in Ruhe die gesamte Darstellung nochmals überblicken und Fragen oder Anmerkungen formulieren.

**Fragerunde.** Fragen während der Präsentation zuzulassen, macht eine Präsentation zwar oft lebendiger. Viele ad hoc gestellte Fragen werden jedoch möglicherweise durch die nachfolgenden Erklärungen ohnehin beantwortet. Die vorbereitete, konzise Struktur einer Präsentation läuft Gefahr, unterlaufen zu werden. Eine gute Alternative ist es, nach einer Präsentation explizit die Möglichkeit vorzusehen, Fragen zu stellen. Alle sollten etwas Zeit bekommen, offene Fragen zunächst aufzuschreiben und sich dabei den Inhalt der Präsentation nochmals zu vergegenwärtigen. Eine wichtige Regel ist, dass sich Zuhörende auf die Formulierung von Fragen beschränken sollten. Noch ist nicht die richtige Zeit für umfassende Erörterungen, Bewertungen und ausschweifende Diskussionen. Zunächst geht es darum, durch Feedback Unklarheiten zu beseitigen.

**Feedback.** Gutes Feedback zu geben, ist nicht trivial. Immer sollten der Respekt und die Anerkennung für die anderen Teilnehmenden erkennbar bleiben – auch und gerade dann, wenn Aspekte kritisiert werden. Es sollte deutlich werden, wenn es um eine subjektive Meinung geht und sich die rückmeldende Person damit keineswegs anmaßt, für alle urteilen zu können. Ich-Botschaften sind hier das Zauberwort: »Ich finde diesen Aspekt nicht passend, zumindest wirkt er auf mich, als gehöre er nicht zum Thema.« Bei allen Vorsichtsmaßnahmen, niemandes Stolzes verletzen zu wollen, sollte unbedingt darauf geachtet werden, ehrliches Feedback zu geben und Kritik nicht aus Rücksicht auszulassen. Rückmeldungen sollten dabei aus zwei Teilen bestehen: dem Urteil und dessen Begründung. Im oben genannten Beispiel wäre es angemessen, das Urteil mit einer Begründung zu hinterlegen: »Ich finde diesen Aspekt nicht passend, weil ...« Auch beim Einholen von Feedback hilft es, darauf zu achten, dass allen Personen Raum dazu gegeben wird, sich einbringen zu können.

**Bewerten.** Das Trennen der eigentlichen Bewertung von der Präsentation, von Fragen und von Feedback gibt Workshops eine klare Struktur. Workshops wirken dadurch weniger komplex, aufkommende Diskussionen werden zielorientiert kanalisiert. Bewertungen können auch durch die Vergabe von Punkten erfolgen: Teilnehmende bekommen eine festgelegte Anzahl zu vergebender Punkte, die auf die unterschiedlichen Optionen aufgeteilt werden sollen. Nicht immer sind extensive Begründungen notwendig – manchmal steht das Ergebnis als Überblick im Vordergrund. Eine interessante Alternative stellt das Bewerten unter Rückgriff auf ein Diagramm dar.

Stellen wir uns beispielsweise vor, es gäbe vier zu bewertende Alternativen. Hier können wir einen Raum aufspannen, dessen Achsen mit für die Aufgabenstellung relevanten Dimensionen – etwa *Kosten* bzw. *Nutzen* – beschriftet sind. Teilnehmende können Alternativen nun relativ zueinander in diesem Raum platzieren, auf der Basis von Diskussionen verschieben und auf diese Weise ein aussagekräftiges, kompaktes und anschauliches Bewertungsdiagramm erstellen.

Im Verlauf dieses Buches werden wir weitere Methoden zur Durchführung von Bewertungen vorstellen. Hierzu gehören das *Plädoyer* (vgl. Kapitel »Workshop: Ideation«) oder auch die *Priorisierungsmatrix* (vgl. Kapitel »Workshop: Roadmap«).

**Sammeln.** Das Sammeln von Aspekten oder Ideen ist ein zentrales Element in der Lösungsfindung: Durch das gezielte Sammeln lässt sich der Lösungsraum in der Breite erkunden. Zur Verdeutlichung dieses Aspekts hilft ein Gedankenexperiment: Zeichnen Sie, nur grob, fünf unterschiedliche Bäume. Danach wählen Sie einen dieser Bäume als Ausgangspunkt und erstellen für eben diesen fünf Varianten. Es wird Ihnen nicht schwerfallen, diesen Prozess nun zu wiederholen. Wieder wählen Sie eine Variante aus und zeichnen fünf weitere Untervarianten, bei denen Sie gezielt Attribute variieren. Wenn Sie den Prozess für alle ursprünglich gezeichneten Bäume wiederholten, so hätten Sie nun bereits 125 Varianten von Bäumen gezeichnet. Sie haben durch Variation von wenigen Bäumen gezielt Vielfalt erstellt – statt eben einfach »nur« 125 ähnliche Bäume zu zeichnen. Beim Sammeln kann es hilfreich sein, sich gegenseitig zu inspirieren. Stellen Sie sich vor, Sie könnten – während Sie die Bäume zeichnen – bei anderen etwas »spionieren«. Dies würde Sie auf weitere Ideen zur Variation bringen und Sie könnten diese wiederum mit Ihren Ideen kombinieren. Einige Ideation-Methoden, die wir in diesem Buch vorstellen – zum Beispiel die Methode 6-3-5 (vgl. Kapitel »Workshop: Ideation«) –, beinhalten genau dieses Element der wechselseitigen Inspiration.

**Gruppieren.** Auf das Sammeln folgt meist ein Gruppieren oder Kategorisieren. Hierbei geht es darum, mehrfach genannte oder bedeutungsähnliche Aspekte in ihrem Kern zusammenzufassen, um hierdurch eine Ordnung zu finden. Die Ordnung kann dabei durchaus vorgegeben sein. Wir haben ein solches Konstrukt zuvor bereits eingesetzt: die (Proto-)Problem Statement Map. Bei dieser sind einzelne Kategorien vorgegeben und wir sammeln nun Annahmen, Beobachtungen oder weitere Aspekte zu diesen Kategorien. Eine kategorisierende Ordnung kann auch erst bei der Betrachtung der unterschiedlich gesammelten Nennungen gefunden werden. Dies ist etwa bei Techniken wie dem Open Card Sorting oder dem Affinity Diagramming der Fall.

Eine weitere Option ist die Zuordnung gesammelter Elemente zu einem zeitlichen Ablauf: Als Beispiel können wir die später in diesem Kapitel diskutierten Journey Maps nennen, in denen wir Findings einzelnen Schritten eines Prozesses zuordnen.

**Ausarbeiten.** Schließlich bleibt noch ein letzter Baustein erfolgreicher Workshops: das Ausarbeiten. Beispiele für das Ausarbeiten sind etwa das Anfertigen einer Skizze, das Beschreiben eines Szenarios oder das Bauen eines Prototyps. Auch wenn wir Ausarbeitungen im Rahmen eines Workshops durchführen, ist es oft hilfreich, dies nicht mit allen Teilnehmenden gleichzeitig, sondern in Kleingruppen oder auch in Form von Einzelarbeiten zu tun.

Die unterschiedlichen, oben dargestellten Schritte können wir, wie in Abbildung 11 dargestellt, zu einem beispielhaften Workshop-Ablauf kombinieren.

Stellen wir uns vor, wir möchten zwischen unterschiedlichen Benutzergruppen differenzieren. Wir könnten zunächst Annahmen zu möglichen Kandidaten aus der Menge der Nutzenden sammeln, auf Karten schreiben und diese anschließend gruppieren. Für eine detaillierte Ordnung, Bewertung und Ausarbeitung der resultierenden Gruppen können wir dann die Teilnehmenden des Workshops in kleinere Teams aufteilen. Anschließend können die Teams einander die aus den Gruppenarbeiten hervorgegangenen Ergebnisse vorstellen und Feedback geben.

## Proto-Personas

Auf den Karten in Abbildung 7 auf Seite 29 finden wir gleich mehrere Annahmen zu Nutzerinnen und Nutzern von 4Service. Tim regt an, die formulierten Annahmen zu konsolidieren: »Wir brauchen möglichst umfassendes und belastbares Wissen über die User von 4Service. Nur so können wir ihre Aufgaben, Arbeitskontexte, Probleme und hieraus resultierende Bedürfnisse verstehen und relevante Nutzungsanforderungen ableiten.« Tim schlägt vor, in der nächsten Workshop-Phase gemeinsam mit Repräsentanten des Produktmanagements *Proto-Personas* zu erstellen.

> **Proto-Persona:** Eine Proto-Persona beschreibt einen hypothetischen »archetypischen« User. Sie steht exemplarisch für eine konkrete Person, die eine Benutzergruppe mit ähnlichen Anforderungen an den Funktionsumfang und das Interaktionsdesign eines Produkts repräsentiert. Eigenschaften einer solchen Proto-Persona stellen zunächst Annahmen dar, die im weiteren Projektverlauf durch Nutzerforschung validiert werden sollten.

Die Idee hinter dem Konzept einer Proto-Persona ist sehr einfach: Auch hier geht es darum, Annahmen explizit zu machen – in unserem Fallbeispiel stehen Annahmen zu den Nutzenden von 4Service im Fokus. Tim möchte wissen, durch welche Merkmale relevante User von 4Service charakterisiert werden können. »Wir beschäftigen uns zunächst nur mit unseren Annahmen zu Nutzenden – über Methoden, die uns dabei helfen, herauszufinden, wie es tatsächlich aussieht, werden wir später sprechen«, sagt Tim.

Tim bittet die Anwesenden darum, ihre individuellen Vorstellungen zu Nutzerinnen und Nutzern von 4Service in Proto-Personas zusammenzufassen. Er schlägt vor, dass dies zunächst alle Anwesenden für sich alleine tun sollen. Die Ergebnisse können nachfolgend einander gegenübergestellt werden, um Gemeinsamkeiten und Unterschiede im Team herauszuarbeiten. Tim greift an dieser Stelle auf eine der zuvor dargestellten Workshop-Regeln zurück: einsam und gemeinsam.

»Nehmt bitte ein Flipchart-Blatt und faltet es einmal entlang der langen Seite: Wenn ihr es wieder aufklappt, habt ihr zwei Hälften. Teilt nun bitte das Blatt entlang der kurzen Seite und ihr erhaltet durch die Faltung vier Quadranten auf dem Blatt.« Tim demonstriert das Falten. »Dies gibt uns eine Struktur zur Beschreibung von Proto-Personas. Damit wir eine Proto-Persona einfach ansprechen können, überlegt euch bitte einen aussagekräftigen Namen und schreibt ihn groß links oben auf das Blatt. Am besten benutzt ihr einen Namen, mit dem deutlich wird, welche Nutzerin damit referenziert wird – ›Marie Beraterin‹ wäre ein solcher sprechender Name. Nun müssen wir nicht mehr vage von ›Nutzerin‹ oder ›Nutzer‹ reden, sondern haben einen Namen, der es uns erlaubt, die Proto-Persona konkret anzusprechen. Wir werden der Persona gleich durch die Nennung relevanter Eigenschaften Leben einhauchen. Schreibt bitte im oberen linken Quadranten der Seite auch auf, welche Ausbildung ihr euch als typisch für die Proto-Persona vorstellt.«

»Darunter könnt ihr dann den Kontext kurz umschreiben, in dem diese Nutzerin arbeitet, und welche typischen Arbeitsaufgaben eure Proto-Persona zu erledigen hat. In einem Segment solltet ihr die Ziele aufführen, von denen ihr glaubt, dass sie für die Proto-Persona bei der Leistungserfassung wichtig sind. Schließlich sollten wir noch die Erwartungen nennen, die die Proto-Persona an ein Modul zur Leistungserfassung nach eurer Ansicht hat. Oft ist es auch hilfreich, zu erwähnen, welche Frustrationspunkte eine Persona bei der Durchführung eines Prozesses hat – oder haben könnte.« Tim gibt den Teilnehmenden etwas Zeit zur inhaltlichen Skizzierung der Proto-Personas und dem Füllen der Felder. Anschließend bittet er sie darum, sich gegenseitig die individuell erstellten Proto-Personas vorzustellen. Abbildung 12 zeigt zwei Proto-Personas als mögliche Ergebnisse der Einzelarbeiten.

Die Vorstellung der angefertigten Proto-Personas hat eine überraschende Vielfalt zum Ergebnis – offensichtlich gehen die am Workshop Teilnehmenden von recht unterschiedlichen Nutzenden aus. Auch wenn oft von »Nutzerin« oder »Nutzer« gesprochen wird, können sich hinter den gleichen Wortmarken doch sehr unterschiedliche Vorstellungen verbergen. Nach der Präsentation und dem erfolgten Feedback konsolidieren die Teilnehmenden die Proto-Personas und halten schließlich die folgenden anschaulichen Namen für ihre fiktiven Proto-Personas fest:

- Peter Projektleiter
- Marc Teamleiter
- Simone Extern
- Marie Beraterin
- Monika Mitarbeiterin-Großprojekt
- Patrick Zentrale Dienste
- Laura Administration
- Lisa Personal

## Theorie: Personas und Proto-Personas

Wir haben im dargestellten Workshop gesehen, wie Teilnehmerinnen und Teilnehmer ihre Vorstellungen zu Usern von 4Service konkretisierten, indem sie ihre unterstellten Eigenschaften in Proto-Personas destillierten. Annahmen über Nutzende werden durch die Formulierung von Proto-Personas explizit gemacht – dennoch bleiben es *Annahmen*. Im Unterschied dazu halten *Personas* empirisch belegte Ergebnisse von Forschungsbemühungen zur Dokumentation von Nutzereigenschaften fest und lassen ein belastbares, in der Realität verankertes Bild von Nutzenden eines Produkts oder Service entstehen. Wie eine Proto-Persona steht eine Persona stellvertretend für eine Nutzergruppe, die die gleichen Anforderungen an ein Produkt oder einen Service stellt: Im Unterschied zu einer Proto-Persona wurden die unterstellten, relevanten Eigenschaften jedoch durch Nutzerforschung auf ihre Gültigkeit hin überprüft. Die durch eine Persona repräsentierten User haben ähnliche Ziele und Erwartungen, ähnliche Fähigkeiten und Erfahrungen und agieren in vergleichbaren Nutzungskontexten.

Wenn wir ein Produkt entwerfen, schließen wir aus den Merkmalen einer modellierten Persona auf erwünschte Eigenschaften dieses Produkts, die die Bedürfnisse der Persona adressieren. Damit es einem Projektteam leichter fällt, Empathie zu Nutzenden zu entwickeln, werden Personas mit persönlichen Attributen wie einem Namen, einem Foto sowie ausgewählten demografischen Eigenschaften ausgestattet. An lediglich schmückenden Attributen oder bloßen demografischen Daten sind wir nicht interessiert – das Erstellen einer Persona ist kein Selbstzweck. So ist beispielsweise anzunehmen, dass die Haarfarbe oder der genaue Geburtsort einer Person keinen bedeutsamen Einfluss auf das Design vieler Lösungen hat, die Nutzungshäufigkeit einer konkreten Funktionalität jedoch sehr wohl.

Je nach Projektphase erfüllen Personas unterschiedliche Zwecke. Sie werden zu Beginn eines Projekts eingesetzt, um typische User eines Systems zu repräsentieren und ihre Wünsche, Bedürfnisse und Frustrationspunkte festzuhalten. Beim Entwurf einer Lösung liefern Personas konkrete Zielvorstellungen für unsere Designaktivitäten und geben dadurch diesen eine Richtung. Ein Projektteam kann bei Entscheidungen im Designprozess immer wieder empathischen Bezug auf die inhaltliche Formulierung einer Persona nehmen und fragen, ob diese Persona ein diskutiertes Feature tatsächlich benötigt oder ob sie diese oder jene Designvariante präferieren würde. Die Persona wird zur anschaulichen Stellvertretung für den intendierten User.

Durch den Rückgriff auf Personas haben wir einen Ausgangspunkt zur Identifikation der Bedürfnisse unterschiedlicher Nutzergruppen, die bei der Produktgestaltung Berücksichtigung finden sollten. Werden deutlich unterschiedliche Personas als ähnlich bedeutsam angesehen, so müssen wir Entscheidungen treffen:

- Ein Lösungskonzept ist konfigurierbar und kann hierdurch an die unterschiedlichen Notwendigkeiten und Präferenzen verschiedener Personas angepasst werden.
- Eine Lösung fokussiert primär auf eine bestimmte Persona und geht in Bezug auf andere, sekundäre Personas Kompromisse ein.
- Die unterschiedlichen Bedürfnisse von Personas rechtfertigen eigene Produktversionen: etwa ein Lösungskonzept, das sich an gelegentliche User richtet, und ein anderes, das die Anforderungen von Power-Usern erfüllt.

Zunächst bleiben jedoch die Ergebnisse der Nutzungsforschung abzuwarten: Wir erleben es in Projekten durchaus, dass initial mehrere Proto-Personas entwickelt wurden, deren Bedürfnisse schließlich in einer Persona kondensiert werden konnten.

## Proto-Journey

> **Proto-Journey:** Eine Proto-Journey beschreibt einen hypothetischen Ablauf zur Erreichung eines der Ziele einer (Proto-)Persona. Eine Proto-Journey betrachtet konkret die aktuelle Situation und modelliert explizit nicht die zukünftige Lösung. Die Proto-Journey umfasst zudem vermutete Probleme und erkannte Chancen. Auch diese stellen zunächst Annahmen dar, die im weiteren Projektverlauf validiert werden sollten.

Während in Proto-Personas die angenommenen Herausforderungen und Erwartungen aus Sicht unterschiedlicher Nutzergruppen beschrieben werden, helfen Proto-Journeys, Handlungen in ihren zeitlichen Abläufen zu ordnen. Die durch Proto-Journeys geschaffene zusätzliche Prozessperspektive hilft, unsere Annahmen zu vervollständigen und das Team auf die nachfolgende Research-Phase vorzubereiten. Der in Proto-Journeys postulierte Ablauf richtet die Aufmerksamkeit in der Nutzerforschung auf zu beobachtende Ereignisse, ihre Sequenz, Häufigkeit und Variation und erlaubt damit eine gezielte Validierung.

Auch bei Proto-Journeys geht es darum, Annahmen explizit zu machen – in unserem Fallbeispiel stehen Annahmen zu den Abläufen bei der Erfassung zu berechnender Leistungen im Fokus. Tim möchte wissen, wie der heutige Workflow bei der Erfassung von Leistungen für eine bestimmte Nutzergruppe aussieht. Er möchte im Workshop den Erfassungsablauf für Marie Beraterin beschreiben.

Tim fragt nach einem Beispiel für eine solche Erfassung, im Team wird die folgende Geschichte erarbeitet:

»Marie Beraterin« erfasst ihre Leistungen eines Tages kurz vor Feierabend. Sie überprüft, welche die letzten von ihr erfassten Leistungen waren, und findet die Einträge vom Vortag. Sie trägt die Leistungen des heutigen Tages ein: In der Regel sind dies etwa zwei bis sechs Einträge. Anschließend validiert sie ihre Einträge, indem sie die Gesamtsumme des Tages betrachtet. Nachdem Marie ihre Leistungen eingetragen hat, erfasst sie ihre Spesen. Dazu sucht sie die Belege zusammen und trägt die Summe ein. Die Belege klebt sie auf ein Papier und legt es in ihre Spesenablage. Am Monatsende überprüft Marie alle Leistungseinträge und bestätigt diese anschließend. Sie weiß, dass auf dieser Grundlage die Projektleitung die Leistungen dann den jeweiligen Kunden in Rechnung stellt.

Tim dokumentiert die wesentlichen Ereignisse der Geschichte auf Karten und ordnet diese in ihrem zeitlichen Ablauf an (Abbildung 13).

Über den skizzierten grundlegenden Ablauf der Leistungserfassung herrscht Konsens unter allen Teilnehmenden des Workshops. Tim fragt nun nach vermuteten Problemen bzw. Chancen zur Optimierung des Ablaufs. Er bittet alle am Workshop teilnehmenden Personen darum, sich zunächst alleine Gedanken zu Chancen und möglichen Barrieren der aktuellen Leistungserfassung zu machen. Alle sollen eigene Karten schreiben und diese anschließend vorstellen. Vermutete Probleme werden dann direkt unter den jeweiligen Schritten platziert. Auf diese Weise entsteht das nachfolgende Bild (Abbildung 14).

Tatsächlich zeigt der Ablauf in der Proto-Journey gegenüber dem Entwurf des Proto-Problem Statement einige zusätzliche Probleme auf. So war im ursprünglichen Proto-Problem Statement nicht erwähnt, dass die Leistungserfassung »bei vielen Projekten unübersichtlich« ist, »keine Übersicht über bereits erfasste Spesen« gegeben ist und »kein einfaches Umbuchen« als Problem angesehen wird. Ebenfalls neu ist die Vermutung, dass Leistungseinträge direkt aus der Wochenplanung übernommen werden können. In einer Wochenplanung wird festgelegt, welche Person mit welchem Aufwand an welchem Projekt arbeitet. Dies könnte eine gute Vorlage liefern, sodass ein geringerer Aufwand für die nachträgliche Erfassung der tatsächlich erbrachten Leistungen anfällt.

Offensichtlich vervollständigt sich das Bild – dem Team wird zunehmend klar, auf welche Aspekte in der Nutzerforschung geachtet werden sollte.

## Annahmen-Map

In den vergangenen Stunden des Workshops wurde eine ganze Reihe von Annahmen durch das Team zusammengetragen. Noch aber ist nicht geklärt, ob und wie bedeutsam oder kritisch die identifizierten Annahmen sind. Die am Workshop Teilnehmenden betrachten daher die verschiedenen Annahmen aus der Erarbeitung des Proto-Problem Statement, der Proto-Personas und der Proto-Journey. Tim bittet darum, darin enthaltene zentrale Annahmen herauszustellen und explizit auf Annahmenkarten zu formulieren (Abbildung 15).

Tim möchte gemeinsam mit dem Team den »Wissensgrad« – das Vertrauen in die Korrektheit der Annahmen – und dessen erwarteter Einfluss auf die Lösungsfindung einschätzen. Er nutzt hierzu ein Diagramm auf Basis eines Koordinatensystems, dessen Achsen er mit *Unsicherheit* und *Einfluss* beschriftet hat: Je weiter rechts eine Person eine Annahmekarte platziert, umso geringer ist das Vertrauen in den Wahrheitsgehalt, je höher sie positioniert wird, umso höher ist ihr erwarteter Einfluss. Die Positionierung der Karte im Koordinatensystem macht den Status für alle Beteiligten deutlich. Rechts im Koordinatensystem stehen entsprechend die Annahmen, bei deren Korrektheit Zweifel bestehen *(Unsicherheit)*. Die Platzierungshöhe einer Annahmekarte zeigt den erwarteten Schaden *(Einfluss)* an, sollte sich diese Annahme später als falsch erweisen. Karten, die einen hohen Einfluss wiedergeben, sind demnach im oberen Teil des Koordinatensystems platziert.

Tim bittet das Team, alle Annahmenkarten im Koordinatensystem (Abbildung 16) zu platzieren und dann gemeinsam über eine Positionierung abzustimmen. Schnell entbrennen heftige Diskussionen: Annahmenkarten werden mit unterstützenden Argumenten immer wieder hin- und herbewegt, bis das Team schließlich mit dem Ergebnis der Annahmen-Map in Abbildung 17 zufrieden ist.

Tim betrachtet das Koordinatensystem und stellt die für das Projekt wesentlichen Annahmen heraus: »Die in der Annahmen-Map oben rechts platzierten Annahmen sind kritisch. Sie haben einen hohen Einfluss – aber wir sind uns über ihren Wahrheitsgehalt nur wenig sicher.«

Das bislang erreichte Ergebnis fasst Tim zusammen: »Alle Annahmen innerhalb des rechten oberen Quadranten sollten wir auf jeden Fall überprüfen, sie sind riskant. Erst auf der dann konsolidierten Grundlage können wir belastbare Entscheidungen treffen.«

Die kritischste Annahme bezieht sich auf die Nutzenden von 4Service. Es wird hypothetisch von insgesamt acht verschiedenen Proto-Personas gesprochen. Zudem liegt die grundlegende Annahme vor, dass Nutzende – wenn es ein besseres Erfassungssystem gäbe – mehr Leistungen erfassen würden. Das heißt, es wird vermutet, dass es wegen Unzulänglichkeiten des aktuellen Erfassungsmoduls zu nicht abgerechneten Leistungen kommt.

Die aufgestellten kritischen Annahmen müssen sorgfältig überprüft werden. Sie sind offenbar hochrelevant für das Produkt – aber das Team weiß zu wenig über ihre Gültigkeit.

## Zusammenfassung Scoping

Den Projektstart konnten wir nun beobachten: Das Projekt wurde von der Auftraggeberin an ein interdisziplinäres Team übergeben. Die Auftraggeberin verfolgt Ziele, deren Erreichung sie an das Team delegiert. Die Arbeit eines Teams kann nur erfolgreich sein, wenn die Ziele eines Auftrags möglichst präzise und umfassend verstanden werden. Bereits in den Zielen kann eine Barriere für den weiteren Projektverlauf liegen – die vorgegebenen Ziele des Auftrags können unangemessen sein. Ebenso könnte von Randbedingungen ausgegangen werden, die nicht (alle) korrekt sind. Die Auftraggeberin könnte also Annahmen getroffen haben, auf die wir uns nicht verlassen können. Gleiches gilt für weitere Stakeholder. Auch diese können Interessen verfolgen, die auf unsicheren Annahmen basieren.

In einem ersten Schritt erkundete das Team daher das Anliegen der Auftraggeberin und weiterer Stakeholder mit dem Ziel, dieses in den relevanten Attributen zu verstehen. Hierzu wurden die bedeutsamen Attribute des Auftrags in einem Proto-Problem Statement geklärt. Um festzustellen, ob das Proto-Problem Statement korrekt und belastbar ist, hinterfragte es das Team durch die Anfertigung einer Annahmen-Map. Gemeinsam mit der Auftraggeberin untersuchte das Team, welche Annahmen dem Auftrag zugrunde liegen und welche Mitarbeitenden das aktuelle System nutzen. Die Annahmen wurden nach dem Grad ihres vermuteten Einflusses auf den Projekterfolg beurteilt. Durch das Vorgehen im Scoping-Workshop wurden der Auftraggeberin und allen Stakeholdern die konkrete Ausgangssituation und deren Rahmenbedingungen vor Augen geführt.

Eine Betrachtung der wichtigen Elemente, bei denen wir unsicher sind, verdeutlicht, welche Klärungen noch notwendig sind: Es liegt nun eine Grundlage zur Planung der Nutzerforschung vor. Schließlich werden es Nutzerinnen und Nutzer sein, die die nächste Version von 4Service als gelungene Unterstützung bewerten – oder das nächste Release ablehnen.
